\documentclass{article}
\usepackage{graphicx} % Required for inserting images

\title{10-30 ECE544 Weekly Report}
\author{Gregory Stokes, Randy Peirce}
\date{October 2026}

\begin{document}

\maketitle

\section{Progress and tasks completed this week:}
\subsection{Gregory Stokes}
    \begin{itemize}
        \item Setup Anaconda and initual packages on computer
        \item Research Netlists and a bit on DRC
        \item Failed Attempt to modify PICBench to use Yaml netlist rather than the sax notaion.
        \item Began creating a converter form sax notation to Yaml
    \end{itemize}
\subsection{Randy Peirce}
    \begin{itemize}
        \item Setup shared Conda/pip environments and initialized the project GitHub repository
        \item Installed and validated core photonics libraries (gdsfactory, ubcpdk, sax, and KLayout)
        \item Created a multi-layer GDSII export smoke test and verified KLayout database parsing
        \item Reviewed baseline photonic benchmark netlists (PICBench, LiDAR)
    \end{itemize}

\section{Plan for Next Week:}
\subsection{Gregory Stokes}
    \begin{itemize}
        \item Finish Test and Verification
        \item Finish Research on DRC
        \item Create Script to DRC on Klayout
        \item Test the script on various Yaml models to verify correctness.
    \end{itemize}

 \subsection{Randy Peirce}
    \begin{itemize}
        \item Implement 3D loss-aware cost function in Python calibrated for 1550 nm (Si vs. SiN)
        \item Define 3D routing state space with vertical taper footprints and U-turn pruning
        \item Implement bounding-box search corridors to ensure graph scalability
        \item Setup 2D baseline routing harness and extract initial benchmark topology (e.g., Benes $4\times4$)
    \end{itemize}


\end{document}
